\section{PIM Design Principles}
\label{sec:principles}
%
Since state update operations are memory-bound, PIM appears to be a
promising solution for acceleration.
%
While prior works have already explored PIM acceleration for
transformer-based LLM serving~\cite{heo2024neupims,
hyun2024pathfinding, newton, hbm_pim}, we observe that SU-LLMs have
significantly different performance characteristics, necessitating
distinct design decisions.
%
In this section, we share our empirical insights and the
corresponding principles that govern our accelerator design.

\subsection{Principle 1: Maximizing Hardware Resource Sharing for
Area Efficiency}
\label{subsec:area_throughput}
%
There are two primary approaches to design PIM for state update
acceleration: (1) time-multiplexed PIM and (2) pipelined PIM.
%
The time-multiplexed design only implements basic multiplication and
addition units to reduce area overhead, similar to HBM-PIM~\cite{hbm_pim}.
%
The pipelined design maximizes throughput, by implementing the entire
sequence of operations in a pipelined manner.
%

%
\niparagraph{Tradeoff between area and throughput.}
%
Figure~\ref{fig:challenges_pim} shows the normalized throughput of
state update operations of SU-LLMs on an A100 GPU and the two PIM
designs during the generation phase at batch size 128.
%
It also presents the respective area overhead of the two PIM designs.
%
We adopt a per-bank PIM design (each bank equipped with its own
processing unit) for both PIM architectures, matching the channel
bandwidth to that of the A100.
%
Detailed PIM configurations and simulator/RTL implementation are
provided in Section~\ref{sec:method}.
%
While the time-multiplexed design offers a modest 2.8$\times$
throughput improvement over the GPU, the pipelined design achieves a
more significant 4.3$\times$ improvement.
%
Meanwhile, the time-multiplexed design has only 17.8\% area overhead,
whereas the pipelined design incurs a much larger area overhead of
32.4\% (>25\%), posing practical deployment challenges.
%
Neither design offers both high throughput and low area overhead.
%

%
\niparagraph{Achieving both high throughput and area efficiency.}
%
We argue that even if the number of processing units in the per-bank
pipelined design is halved, the same throughput can still be
maintained, thereby achieving both high throughput and area efficiency.
%
Even per-bank pipelined designs cannot fully utilize each processing
unit because state updates require both read and write, and row
buffers cannot perform both simultaneously.
%
During writes, no input is supplied to the processing unit.
%
With judicious dataflow design, two banks can share a single
processing unit, allowing continuous input from both banks without
throughput loss.
%
We call this technique \emph{access interleaving} and detail it in
Section~\ref{subsec:access_interleaved_iteration_parallelism}.

%

\subsection{Principle 2: Achieving Both Accuracy and Efficiency from
Low-Precision Arithmetic}
\label{subsec:area_accuracy}
%
As discussed in Section~\ref{subsec:challenges_quantization},
quantizing the state is another promising approach to alleviate the
pressure on memory bandwidth during state update operations.
%
To leverage quantization in PIM, we conduct experiments to evaluate
the tradeoff between area and accuracy aiming to achieve the best of both.
%

%
\niparagraph{Tradeoff between area and accuracy.}
%
Figure~\ref{fig:quantization_pareto} illustrates the area-accuracy
tradeoff space for various low-precision formats, evaluated using the
Mamba-2 on the Wikitext-2~\cite{merity2016pointer} dataset with a
per-bank pipelined PIM design.
%
Detailed RTL implementation methodology is provided in Section~\ref{sec:method}.
%
As described in Section~\ref{subsec:challenges_quantization},
\texttt{fp8} formats suffer from a significant increase in
perplexity, while \texttt{int8} and \texttt{MX8} achieve perplexity
levels comparable to \texttt{fp16}.
%

%
However, our analysis suggests that \texttt{int8} incurs substantial
area overhead, while \texttt{MX8} is significantly more area-efficient.
%
The high area overhead of \texttt{int8} stems from the need for
element-wise addition during state updates, which the scaling-based
integer format cannot directly handle.
%
This necessitates dequantization (multiplying each integer element by
the scaling factor) and re-quantization (normalizing elements based
on the max value), requiring additional arithmetic units and comparison logic.
%
Conversely, \texttt{MX} simplifies operand alignment for addition by
sharing the exponent bit within a group, enabling direct operations
through simple shifting without dequantization, the details of which
will be discussed in Section~\ref{subsec:microarchitecture}.
%

Additionally, the results show that stochastic rounding imposes
minimal area overhead.
%
It requires a Linear Feedback Shift Register (LFSR) for random number
generation and a simple addition unit to add these numbers to the
mantissa, both of which are area-efficient~\cite{9773221}.
%
We conclude that among the quantized formats, \texttt{MX8} with
stochastic rounding emerges as a Pareto-optimal choice, offering
superior area efficiency while maintaining high accuracy.
